\chapter{Dérivés sur l'inflation}\label{ch:rc:inflation-derivatives}

En octobre 2022, l'indice des prix de détail du Royaume-Uni se situait 14.2~\% au-dessus de son
niveau un an plus tôt. De nombreux fonds de pension britanniques ne promettent pas à leurs membres
une indexation intégrale de l'inflation~: leurs revalorisations sont égales à la variation annuelle de l'indice, plancher à zéro et
plafonnée à 5~\% (2.5~\% pour les droits acquis depuis 2005). Les fonds qui avaient couvert ces
engagements au moyen de simples swaps sur indice détenaient des couvertures qui versaient la totalité des
14.2~\%, alors que leurs passifs ne devaient augmenter que de 5~\%~; l'écart entre un indice et un indice plafonné,
réinitialisé chaque année, était devenu le chiffre le plus important de leur rapport de risque. Tout
produit d'inflation au-delà du swap zéro-coupon du livre~2, chapitre~11,
dépend de la manière dont l'indice évolue d'une année à l'autre, et c'est une question qui relève
d'un modèle. Ce chapitre construit le modèle standard par analogie avec le change,
calcule la convexité qui distingue les produits year-on-year des produits zéro-coupon,
ajoute la saisonnalité, et price l'\omterm{def:rc:inflation-derivatives:lpi}{indexation plafonnée} (limited price indexation, LPI) par simulation.

\section{L'analogie avec les devises}

Traitons la monnaie nominale comme une devise et la monnaie «~réelle~», c'est-à-dire une unité du panier,
comme une autre. Une obligation nominale verse de la devise~; une obligation indexée verse un panier, valant
$I(T)$ unités de devise. L'indice joue le rôle de taux de change entre les deux, les facteurs
d'actualisation réels $P_R(t,T)$ sont les facteurs d'actualisation de la devise étrangère, et
l'indice à terme est un taux de change à terme~:
\[
I(t,T) = I(t)\,\frac{P_R(t,T)}{P_N(t,T)},
\]
une martingale sous la probabilité forward nominale en $T$. Un swap d'inflation zéro-coupon
fixe $(1+k_T)^T = I(0,T)/I(0)$, de sorte que la courbe réelle d'aujourd'hui se lit à partir des swaps~:
$P_R(0,T) = P_N(0,T)(1+k_T)^T$.

\begin{definition}[Modèle de Jarrow--Yildirim]\label{def:rc:inflation-derivatives:jy}
Le \emph{modèle de Jarrow--Yildirim}\index{modèle de Jarrow--Yildirim} (2003) modélise le
taux court nominal et le taux court réel comme deux processus de Hull--White, et
l'indice comme un taux de change lognormal entre eux~:
$dI/I = (n_t-r^R_t)\,dt+\sigma_I\,dW_I$ sous la probabilité nominale, avec des corrélations
entre les trois facteurs d'aléa. Comme pour un taux court étranger vu depuis le pays d'origine
(l'ajustement quanto du livre~5, chapitre~17), le taux réel acquiert la dérive
$-\rho\,\sigma_R\sigma_I$ sous la probabilité nominale, $\rho$ étant la corrélation du
taux réel avec l'indice.
\end{definition}

Le chapitre maintient les taux nominaux déterministes, ce qui isole les effets du
taux réel et de l'indice~; le modèle complet ajoute des termes en volatilité nominale et
ses corrélations de même forme.

\begin{omfigure}
\begin{tikzpicture}[font=\small,
  box/.style={draw,thick,rounded corners=2pt,align=center,minimum height=1cm,minimum width=3.4cm,inner sep=4pt},
  arr/.style={-{Stealth[length=2.5mm]},thick}]
\node[box,draw=omDef,fill=omDef!8] (n) at (0,0) {économie nominale\\taux $n_t$, obligations $P_N$};
\node[box,draw=omThm,fill=omThm!8] (r) at (8,0) {économie réelle\\taux $r^R_t$, obligations $P_R$};
\draw[arr,omProp] (r.160) -- node[above=4pt,align=center] {indice $I(t)$~: monnaie par panier\\(le taux de change)} (n.20);
\node[align=center,font=\footnotesize] at (4,-1.35) {monnaie domestique $\leftrightarrow$ nominale~;\quad devise étrangère $\leftrightarrow$ réelle~;\quad change à terme $\leftrightarrow$ indice à terme};
\end{tikzpicture}

\omcaption{L'analogie des devises qui sous-tend les modèles d'inflation~: la monnaie réelle est une
devise étrangère dont le taux de change est l'indice des prix, de sorte que la mécanique
du pricing inter-devises, y compris les dérives quanto, se transpose.}\label{fig:rc:inflation-derivatives:analogy}
\end{omfigure}

\section{Zéro-coupon contre year-on-year}

\begin{definition}[Swap d'inflation year-on-year]\label{def:rc:inflation-derivatives:yoy}
Un \emph{swap d'inflation year-on-year}\index{swap d'inflation year-on-year} échange,
chaque année $i$, la variation annuelle de l'indice $I(T_i)/I(T_{i-1})-1$ contre un taux fixe,
sur un notionnel qui ne croît pas avec l'indice.
\end{definition}

Un swap zéro-coupon verse un seul rapport $I(T)/I(0)$ en $T$~: sa valeur ne requiert que
l'indice à terme. Une jambe year-on-year verse un rapport de deux valeurs futures de l'indice, payé
à la date la plus tardive~; sa valeur requiert l'espérance du rapport sous la probabilité
nominale de la date de paiement, et le numérateur et le dénominateur du rapport sont corrélés
par le taux réel.

\begin{definition}[Ajustement de convexité year-on-year]\label{def:rc:inflation-derivatives:yoyadj}
L'\emph{ajustement de convexité year-on-year}\index{ajustement de convexité year-on-year}
de la période $i$ est l'écart entre $\E^{T_i}_N[I(T_i)/I(T_{i-1})]$ et le rapport des
indices à terme $I(0,T_i)/I(0,T_{i-1})$.
\end{definition}

\begin{proposition}[Espérance year-on-year]\label{prop:rc:inflation-derivatives:yoy}
Dans le \omterm{def:rc:inflation-derivatives:jy}{modèle de Jarrow--Yildirim} à taux nominaux déterministes, le taux réel
$r^R = f_R(0,t)+x_t$ avec les paramètres de Hull--White $(\kappa,\sigma_R)$ et une corrélation $\rho$
entre le taux réel et l'indice,
\[
\E_N\Bigl[\frac{I(S)}{I(T)}\Bigr] = \frac{I(0,S)}{I(0,T)}\,e^{C(T,S)},\quad
C(T,S) = -B(T,S)\Bigl(\frac{\sigma_R^2}{2\kappa^2}\bigl(1-e^{-\kappa T}\bigr)^2-\rho\,\sigma_R\sigma_I\,B(0,T)\Bigr).
\]
\end{proposition}
\begin{proof}
En conditionnant en $T$, on obtient $\E_T[I(S)] = I(T)\,P_R(T,S)/P_N(T,S)$, de sorte que
l'espérance du rapport vaut $\E_N[P_R(T,S)]/P_N(T,S)$. Avec
\[
P_R(T,S) = \frac{P_R(0,S)}{P_R(0,T)}\,e^{-Bx_T-\frac12B^2y(T)}
\]
et $x_T$ normale de variance $y(T)$ et de moyenne sous la probabilité nominale $\mu(T)-\rho\sigma_R\sigma_IB(0,T)$,
où $\mu(T) = \int_0^Te^{-\kappa(T-u)}y(u)\,du = \frac{\sigma_R^2}{2\kappa^2}(1-e^{-\kappa T})^2$, il reste à calculer
$\E[e^{-Bx_T}]$.
\end{proof}

\omcode{code/firm/inflopt/firm_inflopt.py}{63}{70}{L'espérance year-on-year et son ajustement de convexité.}\label{lst:rc:inflation-derivatives:yoy}

\begin{example}[Convexité year-on-year en livre sterling]\label{ex:rc:inflation-derivatives:conv}
Sur des courbes illustratives en livre sterling (taux zéro SONIA de 3.6 à 4.2~\%, swaps zéro-coupon RPI
de 3.2 à 3.4~\%) avec $\kappa=5~\%$, $\sigma_R=80$ points de base, $\sigma_I=1.5~\%$ et
$\rho=0.2$, l'ajustement vaut zéro pour la première année, $-1.8$ points de base pour la quatrième,
$-15.2$ pour la dixième (taux à terme 3.42~\%, espérance year-on-year de 3.27~\%) et $-42.5$ pour la
dix-neuvième (\cref{fig:rc:inflation-derivatives:convexity}). Le taux du swap year-on-year à dix ans
vaut 3.22~\%, contre 3.28~\% pour le taux zéro-coupon à dix ans. Une simulation du modèle
reproduit l'espérance de la dixième année, 1.03268, à 1.03277 avec une erreur type de
0.00013.
\end{example}

\begin{omfigure}
\begin{tikzpicture}
\begin{axis}[width=11.5cm,height=5.4cm,
  xlabel={année $i$ de la variation annuelle},ylabel={taux (\%)},
  tick label style={font=\small},label style={font=\small},
  xmin=0.5,xmax=20.5,ymin=2.8,ymax=3.6,grid=major,grid style={gray!25},
  legend columns=2,legend style={font=\footnotesize,at={(0.5,-0.25)},anchor=north,draw=none,fill=none,/tikz/every even column/.append style={column sep=8pt}},
  table/col sep=comma]
\addplot[omDef,thick,mark=*,mark size=1.4pt] table[x=i,y=fwd]{figdata/rates-credit-risk/11-inflation-derivatives/convexity.csv};
\addplot[omThm,very thick,mark=square*,mark size=1.4pt] table[x=i,y=yoy]{figdata/rates-credit-risk/11-inflation-derivatives/convexity.csv};
\legend{ratio d'indices à terme,espérance year-on-year}
\end{axis}
\end{tikzpicture}

\omcaption{Inflation annuelle implicite des swaps zéro-coupon (ratio d'indices à terme)
et espérance du modèle pour la variation de chaque année, que verse un swap year-on-year.
L'écart, soit l'ajustement de convexité, croît avec le début de l'année. Données~:
courbes illustratives en livre sterling, paramètres de Jarrow--Yildirim
de l'\cref{ex:rc:inflation-derivatives:conv}~; le
tutoriel du chapitre.}\label{fig:rc:inflation-derivatives:convexity}
\end{omfigure}

Le signe et la taille dépendent de la corrélation~: à la dixième année, l'ajustement vaut
$-20.6$, $-16.4$ et $-12.2$ points de base pour $\rho = -0.5$, $0$ et $0.5$. Sans volatilité
des taux réels, il s'annule quelle que soit $\rho$~; la seule volatilité de l'indice ne le crée pas.

\section{Saisonnalité dans le modèle}

Les indices de prix ne sont pas corrigés des variations saisonnières, et leurs variations mensuelles suivent le
calendrier~: soldes, tarifs de l'énergie, jours fériés (livre~2, chapitre~11). Un
swap d'inflation ou une obligation référence l'indice d'un mois donné, si bien que l'indice à terme doit
reproduire le profil saisonnier entre les points annuels tout en les respectant.

\begin{method}[Indice à terme saisonnier]\label{met:rc:inflation-derivatives:season}
Estimer l'écart moyen de la variation logarithmique de chaque mois calendaire par rapport
à la moyenne de son année (les facteurs saisonniers, de somme nulle)~; construire l'indice à terme mois par mois
comme l'indice à terme lissé multiplié par les facteurs cumulés. Les points annuels sont inchangés, la
trajectoire entre eux s'infléchit, et tout produit fixant sur un mois autre que les mois de référence des
swaps voit son indice varier.
\end{method}

\begin{example}[Le profil américain]\label{ex:rc:inflation-derivatives:season}
Sur l'IPC américain (non corrigé des variations saisonnières) sur les années complètes 2010--2025, la
variation logarithmique mensuelle dépasse de 0.18, 0.26 et 0.31 points de pourcentage la moyenne de son année, en
janvier, février et mars, et lui est inférieure de 0.36 et 0.35 points de pourcentage en novembre et
décembre (\cref{fig:rc:inflation-derivatives:seasonal}). Ajouté à un indice à terme
lissé à 3~\% par an, ce profil le relève de 1.0 point à fin mai et
de 1.1 point à fin juin, et le ramène au chemin lisse à chaque fin d'année.
\end{example}

\begin{omfigure}
\begin{tikzpicture}
\begin{axis}[width=11.5cm,height=5.2cm,
  xlabel={mois depuis aujourd'hui (janvier)},ylabel={indice à terme (base 100)},
  tick label style={font=\small},label style={font=\small},
  xmin=0,xmax=24,ymin=99.5,ymax=106.5,xtick={0,3,6,9,12,15,18,21,24},grid=major,grid style={gray!25},
  legend columns=2,legend style={font=\footnotesize,at={(0.5,-0.25)},anchor=north,draw=none,fill=none,/tikz/every even column/.append style={column sep=8pt}},
  table/col sep=comma]
\addplot[omDef,thick,dashed] table[x=m,y=smooth]{figdata/rates-credit-risk/11-inflation-derivatives/seasonal.csv};
\addplot[omThm,very thick,mark=*,mark size=1.2pt] table[x=m,y=seasonal]{figdata/rates-credit-risk/11-inflation-derivatives/seasonal.csv};
\legend{lissé à 3~\% par an, avec le profil saisonnier américain}
\end{axis}
\end{tikzpicture}

\omcaption{Indice à terme sur deux ans à 3~\% par an, lisse, et avec les facteurs saisonniers mensuels
de l'IPC américain (années complètes 2010--2025). Les deux passent par les mêmes
points annuels. Source~: US Bureau of Labor Statistics via FRED (série CPIAUCNS, fichier de données du
livre~2)~; le
tutoriel du chapitre.}\label{fig:rc:inflation-derivatives:seasonal}
\end{omfigure}

\section{Options d'inflation et indexation plafonnée}

\begin{definition}[Cap d'inflation]\label{def:rc:inflation-derivatives:cap}
Un \emph{cap d'inflation}\index{cap d'inflation} est une bande d'options versant chacune
$\max(I(T_i)/I(T_{i-1})-1-K,0)$ à $T_i$~: un cap sur l'inflation annuelle (year-on-year), avec des floors
(le plancher de déflation du livre~2, chapitre~11, pour le cas zéro-coupon) définis
de la même façon.
\end{definition}

Dans le modèle, le ratio annuel est lognormal~: sa variance logarithmique est la variance de l'indice
sur l'année, plus la variance de l'intégrale du taux réel de cette année, moins deux fois leur
covariance, et le caplet est la formule de Black appliquée à l'espérance year-on-year.

\begin{example}[Un cap et un floor à dix ans]\label{ex:rc:inflation-derivatives:caps}
Sur les courbes en livre sterling, un cap year-on-year à 5~\% sur les années deux à dix coûte 181.1
points de base du notionnel~; un floor à 0~\% coûte 45.8.
\end{example}

\begin{definition}[Indexation plafonnée]\label{def:rc:inflation-derivatives:lpi}
L'\emph{indexation plafonnée}\index{indexation plafonnée} (limited price indexation, LPI) revalorise un montant
chaque année de la variation annuelle de l'indice, minorée à $f$ et plafonnée à $c$, en capitalisant
le résultat~: après $N$ années, il a crû de $\prod_{i=1}^N\bigl(1+\min(\max(I_i/I_{i-1}-1,f),c)\bigr)$.
Les revalorisations des retraites britanniques suivent ce mécanisme par la loi, avec $[0~\%,5~\%]$ pour les droits acquis avant le 6
avril 2005 et $[0~\%,2.5~\%]$ après.
\end{definition}

La capitalisation de variations plafonnées et minorées dépend du chemin suivi~: le paiement n'est pas une somme de
caplets. Il se valorise par simulation conjointe des ratios year-on-year.

\omcode{code/firm/inflopt/firm_inflopt.py}{119}{127}{La jambe d'indexation plafonnée par simulation des ratios annuels d'indice.}\label{lst:rc:inflation-derivatives:lpi}

\begin{omfigure}
\begin{tikzpicture}
\begin{axis}[width=11.5cm,height=5.2cm,
  xlabel={facteur de croissance après vingt ans},ylabel={part des trajectoires (\%)},
  tick label style={font=\small},label style={font=\small},
  xmin=1.0,xmax=3.1,ymin=0,ymax=5.5,xtick={1,1.5,2,2.5,3},grid=major,grid style={gray!25},
  legend columns=2,legend style={font=\footnotesize,at={(0.5,-0.3)},anchor=north,draw=none,fill=none,/tikz/every even column/.append style={column sep=8pt}},
  table/col sep=comma]
\addplot[omDef,thick,const plot mark mid] table[x=x,y=full]{figdata/rates-credit-risk/11-inflation-derivatives/lpihist.csv};
\addplot[omThm,very thick,const plot mark mid] table[x=x,y=lpi]{figdata/rates-credit-risk/11-inflation-derivatives/lpihist.csv};
\legend{ratio d'indice $I(20)/I(0)$,{LPI $[0~\%,5~\%]$}}
\end{axis}
\end{tikzpicture}

\omcaption{Distribution de la croissance sur vingt ans d'un indice sans plafond et de sa version
$[0\%,5\%]$ à \omterm{def:rc:inflation-derivatives:lpi}{indexation plafonnée} (LPI), 40\,000 trajectoires simulées du modèle sterling~:
le cap supprime la queue droite, le floor les déflations rares. Données~: le
tutoriel du chapitre.}\label{fig:rc:inflation-derivatives:lpi}
\end{omfigure}

\begin{dated}{2026-09}{L'inflation au Royaume-Uni en 2022 et après}\label{dat:rc:inflation-derivatives:uk}
La variation annuelle du RPI fut de 14.2~\% en octobre 2022 (CPI 11.1~\%, CPIH 9.6~\%). Le
gouvernement et l'UK Statistics Authority décidèrent en novembre 2020 que le RPI serait
aligné sur les méthodes et les données du CPIH à partir de février 2030, après la dernière
gilt indexée qui aurait sinon été touchée, ce qui repoussa la courbe des swaps de RPI
au-delà de 2030. Les hausses légales des pensions restèrent plafonnées à 5~\% (droits antérieurs à 2005) ou
à 2.5~\%.
\end{dated}

\section{Tutoriel~: produits year-on-year}\label{tut:rc:inflation-derivatives:tutorial}

\begin{tutorial}
\textbf{Objectif.} Construire les courbes nominale et réelle, calculer la convexité year-on-year,
coter des caps et une jambe LPI, et intégrer la saisonnalité dans un indice forward. \textbf{État final~:}
\cref{fig:rc:inflation-derivatives:convexity,fig:rc:inflation-derivatives:seasonal,fig:rc:inflation-derivatives:lpi}.
\begin{tutsteps}
\item \textbf{Courbes}~: \texttt{real\_curve(nominal, years, zc\_rates)} à partir des
swaps d'inflation zéro-coupon du RPI.
\item \textbf{Convexité}~: \texttt{convexity\_table()}, \texttt{rho\_effect()} et
\texttt{mc\_check()}.
\item \textbf{Options}~: \texttt{caps()} et \texttt{lpi()}.
\item \textbf{Saisonnalité}~: \texttt{us\_seasonals()} avec le \texttt{seasonal\_factors}
du livre~2~; \texttt{fig\_rc\_inflation.py} produit les graphiques.
\end{tutsteps}
\textbf{Pour aller plus loin.} Coter la jambe LPI avec un cap à 2.5~\% et un floor de
$-100~\%$ (sans floor), puis répartir la valeur entre le cap et le floor~; rendre le
taux nominal stochastique dans la simulation et mesurer la variation du swap year-on-year
taux.
\end{tutorial}

\section{Construction~: options d'inflation}\label{bld:rc:inflation-derivatives:inflopt}

\begin{build}
\bfield{Rôle} Le portefeuille d'inflation de l'entreprise miniature~: \omterm{def:rc:inflation-derivatives:yoy}{swaps d'inflation
year-on-year}, caps et floors, et jambes LPI des couvertures de régimes de retraite.
\bfield{Interface} \texttt{LogLinearCurve}, \texttt{real\_curve}~;
\texttt{JYDet(nominal, real, kappa, sigma\_r, sigma\_i, rho)} avec \texttt{yoy\_ratio},
\texttt{convexity}, \texttt{yoy\_log\_variance}, \texttt{yoy\_caplet},
\texttt{yoy\_swap\_rate}, \texttt{zc\_swap\_rate}, \texttt{simulate\_ratios},
\texttt{lpi\_leg}.
\bfield{Règles} Taux nominaux déterministes~; la courbe réelle reproduit exactement les
swaps zéro-coupon~; simulation sous la probabilité nominale avec la dérive quanto~; le
\texttt{firm\_breakeven} du livre~2 pour les facteurs saisonniers.
\bfield{Tests d'acceptation} \texttt{code/firm/inflopt/tests/}~: le modèle reproduit le taux du
swap zéro-coupon et n'a pas d'ajustement la première année ni sans volatilité du taux réel~;
l'espérance year-on-year analytique coïncide avec la simulation~; cap moins floor vaut le
forward~; une jambe LPI aux bornes infinies vaut l'indice, et l'indice simulé coïncide avec
le facteur d'actualisation réel.
\bfield{Pour aller plus loin} Taux nominaux stochastiques (le modèle complet)~; un smile pour
les caplets year-on-year~; saisonnalité dans l'indice forward du modèle aux fixings mensuels.
\end{build}

\begin{omsources}
\item R.\ Jarrow et Y.\ Yildirim, ``Pricing Treasury Inflation Protected Securities
and related derivatives using an HJM model'', \emph{Journal of Financial and
Quantitative Analysis} 38(2), 2003.
\item F.\ Mercurio, ``Pricing inflation-indexed derivatives'', \emph{Quantitative
Finance} 5(3), 2005.
\item ONS, \emph{Consumer price inflation, UK: October 2022}~; HM Treasury et UK
Statistics Authority, réponse sur la réforme du RPI, novembre 2020~; Pensions Act 1995,
section 51.
\end{omsources}

\section{Exercices}

\begin{exercise}[$\star$]\label{exo:rc:inflation-derivatives:1}
Un swap zéro-coupon de RPI à dix ans est à 3.28~\% et le facteur d'actualisation nominal à dix ans
vaut 0.6771. Calculer le facteur d'actualisation réel à dix ans et le ratio d'indice forward.
\end{exercise}

\begin{exercise}[$\star$]\label{exo:rc:inflation-derivatives:2}
Pourquoi n'y a-t-il pas d'ajustement de convexité pour la première année d'un swap year-on-year~?
\end{exercise}

\begin{exercise}[$\star$]\label{exo:rc:inflation-derivatives:3}
Une hausse de pension avec LPI $[0~\%,5~\%]$ suit un indice qui a progressé de 14.2~\%. De combien
la pension augmente-t-elle~? Et avec $[0~\%,2.5~\%]$~?
\end{exercise}

\begin{exercise}[$\star\star$]\label{exo:rc:inflation-derivatives:4}
Expliquer le signe de l'ajustement dans l'\cref{ex:rc:inflation-derivatives:conv} et pourquoi
il devient moins négatif à mesure que $\rho$ augmente.
\end{exercise}

\begin{exercise}[$\star\star$]\label{exo:rc:inflation-derivatives:5}
Lire la \cref{fig:rc:inflation-derivatives:seasonal}~: un swap fait référence à l'indice de mai et un
autre à l'indice de décembre de l'année suivante. Lequel a le forward le plus élevé, par rapport
à la trajectoire lissée, et pourquoi~?
\end{exercise}

\begin{exercise}[$\star\star$]\label{exo:rc:inflation-derivatives:6}
Pourquoi une jambe LPI n'est-elle pas un strip de caps et de floors year-on-year~?
\end{exercise}

\begin{exercise}[$\star\star\star$]\label{exo:rc:inflation-derivatives:7}
\emph{Programmation.} Avec des nombres aléatoires communs, mesurer la variation de la jambe LPI
à vingt ans $[0~\%,5~\%]$ lorsque $\sigma_I$ passe de 1.5~\% à 2.5~\%, et lorsque $\sigma_R$ passe de 0.80~\% à
0.90~\%.
\end{exercise}

\begin{exercise}[$\star\star\star$]\label{exo:rc:inflation-derivatives:8}
\emph{Trouver la faille.} «~Le taux du swap year-on-year est égal au taux du swap zéro-coupon,
puisque les deux moyennent la même inflation forward.~»
\end{exercise}

\section{Problème~: le swap LPI du régime de retraite}

\begin{problem}[{Problème du week-end --- couverture des hausses de pensions plafonnées}]\label{pb:rc:inflation-derivatives:1}
Les engagements d'un régime de retraite augmentent chaque année selon la RPI, avec une LPI $[0~\%,5~\%]$. Son conseiller
propose un swap à vingt ans qui verse, à l'échéance de la vingtième année, la croissance de la LPI d'un notionnel, contre
l'indice non plafonné. Utiliser les courbes en livres sterling et les paramètres du chapitre, 40\,000 trajectoires.

\medskip\noindent\textbf{Partie I --- Les jambes.}
\begin{enumerate}
\item Donner la valeur d'aujourd'hui de la jambe indexée sur l'indice non plafonné, par unité, puis sa valeur exacte.
\item Donner la valeur de la jambe LPI $[0~\%,5~\%]$.
\item Donner la valeur d'une jambe LPI $[0~\%,2.5~\%]$.
\item Pourquoi la jambe LPI est-elle moins chère, et laquelle de ses deux options domine-t-elle~?
\item Que signifie la différence pour un régime couvert par des swaps non plafonnés~?
\end{enumerate}

\medskip\noindent\textbf{Partie II --- Sensibilités.}
\begin{enumerate}[resume]
\item Donner la variation de la jambe LPI lorsque $\sigma_I$ augmente d'un point.
\item Donner la variation lorsque $\sigma_R$ augmente de dix points de base.
\item Après le swap, le régime est-il long ou court en volatilité~?
\item Pourquoi les sensibilités sont-elles calculées sur des nombres aléatoires communs~?
\item Quels instruments de marché permettraient de les couvrir~?
\end{enumerate}

\medskip\noindent\textbf{Partie III --- 2022.}
\begin{enumerate}[resume]
\item L'indice ayant progressé de 14.2~\%, quelle hausse les engagements $[0~\%,5~\%]$ ont-ils reçue~?
\item Que versait une couverture non plafonnée la même année~?
\item Pourquoi 2022 a-t-elle transformé la couverture du régime en position spéculative~?
\item Que change la réforme de la RPI de 2030 pour la jambe LPI~?
\item Comment le problème changerait-il pour un indice assorti d'un plancher de $-100~\%$~?
\end{enumerate}

\medskip\noindent\textbf{Partie IV --- Jugement.}
\begin{enumerate}[resume]
\item Qui peut prendre la contrepartie des swaps LPI~?
\item Pourquoi la corrélation $\rho$ du modèle est-elle un risque que le régime ne voit pas dans les cotations~?
\item Que faudrait-il rapporter aux administrateurs du fonds~?
\item Énoncer le \emph{résultat nommé}~: les valeurs des jambes non plafonnée et LPI, et la
sensibilité de la jambe LPI à la volatilité de l'indice.
\item En une phrase~: que fait un plafond sur les variations annuelles à une promesse de vingt ans~?
\end{enumerate}
\end{problem}

\section{Questions d'entretien}

\begin{interviewq}[$\star$ \iqroles{chercheur, banque}]\label{iq:rc:inflation-derivatives:1}
Expliquez l'analogie avec le change pour l'inflation.
\end{interviewq}

\begin{interviewq}[$\star\star$ \iqroles{chercheur}]\label{iq:rc:inflation-derivatives:2}
Pourquoi un swap year-on-year n'est-il pas valorisé à partir de la seule courbe zéro-coupon~?
\end{interviewq}

\begin{interviewq}[$\star\star$ \iqroles{trader}]\label{iq:rc:inflation-derivatives:3}
Comment la saisonnalité entre-t-elle dans le prix d'un swap d'inflation, et où la trouvez-vous~?
\end{interviewq}

\begin{interviewq}[$\star\star$ \iqroles{trader, banque}]\label{iq:rc:inflation-derivatives:4}
Qu'est-ce que l'\omterm{def:rc:inflation-derivatives:lpi}{indexation plafonnée} sur les prix, et pourquoi rend-elle la couverture des pensions difficile~?
\end{interviewq}

\begin{interviewq}[$\star\star\star$ \iqroles{chercheur, risque}]\label{iq:rc:inflation-derivatives:5}
Quels paramètres du \omterm{def:rc:inflation-derivatives:jy}{modèle de Jarrow--Yildirim} sont observables sur le marché, et lesquels
ne le sont pas~?
\end{interviewq}

\begin{interviewq}[$\star\star\star$ \iqroles{développeur}]\label{iq:rc:inflation-derivatives:6}
Votre indice simulé ne reprice pas les swaps zéro-coupon. Que vérifiez-vous~?
\end{interviewq}
